---
title: "Beamer Presentation (BIO 381 - HW 3)"
author: "Adrian Wiegman"
date: "January 31, 2018"
output:
  ioslides_presentation: default
  beamer_presentation: default
---

```{r setup, include=FALSE}
knitr::opts_chunk$set(echo = FALSE)
```

## Primary objective of the Presentation

- demonstrate use of sequential and simultaneous bullet points
- use colored text
- use an images 
- use an image with a caption
- use an image with an overlayed caption

## Secondary objective of the presentation

- talk about my what i had for breakfast!

## Ingredients of my breakfast 

>- \color{white}e\color{yellow}gg\color{white}s 
>- \color{pink} ham
>- \color{green} green dye
>- \color{cream} butter

## An R vector of what I had for breakfast

```{r cars, echo = TRUE}
Breakfast <- c('green','eggs','and','ham')
print(Breakfast)
```

## An equation for the weight I gained after eating my breakfast

$B = E*n~e~*k~ce~ + GD + H**k~ch~ + B*k~cb~$


- B is the mass of my breakfast (grams), E the mass of on large chicken egg (grams), n~e~ is the number of eggs, k~ce~ is the moisture loss coefficient (g g^-1^) for pan frying one raw egg from a temp of 5degC to a mean tempurature of 70degC 

- GD is the mass of green dye added to the egg yolk prior to frying

- H the mass of a standard 1 cm thick slice of ham shank with a diameter of 15cm (grams), and k~ch~ is the moisture loss coefficient for pan frying a standard slice smoked ham shank up from 5decC to to 70decG

- B the mass of butter added, to greese the pan, k~cb~ is the moisture loss coefficient to bring the eggs and ham from a temp of 5degC to 70degC

## A picture of what I had for breakfast
 \begin{picture}(320,250)
 \put(85,75){\includegraphics[height=6cm]
 {samiam.jpg}}
 \put(00,60){\begin{minipage}[t]{1\linewidth}
 {Choose a point on the unit circle. Connect it to 
 the origin with a line of 
 length one, and denote the angle between that line 
 and the horizontal 
 coordinate axis by $\theta$.}
 \end{minipage}}
 \end{picture}



